\begin{trapbox}
	\begin{description}[style=nextline, leftmargin=2.8em, labelsep=0.5em, itemsep=0.35em]
		\item[\textbf{\textcolor{CTrap}{易错点一（混淆半径公式）}}]
		      把正方体外接球半径误记为
		      \[
			      R=\frac{\sqrt{2}}{2}a
		      \]
		      或
		      \[
			      R=\frac{\sqrt{6}}{2}a.
		      \]

		\item[\textbf{\textcolor{CTrap}{正确理解（正确公式）：}}]
		      正确公式是
		      \[
			      R=\frac{\sqrt{3}}{2}a.
		      \]
		      可通过体对角线
		      \[
			      d=\sqrt{3}a
		      \]
		      且
		      \[
			      R=\frac{d}{2}
		      \]
		      验证。

		\item[\textbf{\textcolor{CTrap}{错因分析（混淆来源）：}}]
		      误以为外接球直径等于面对角线 $\sqrt{2}a$，或者把平面正方形外接圆半径 $\frac{\sqrt{2}}{2}a$ 错套到空间正方体外接球上。
	\end{description}

	\medskip
	\begin{description}[style=nextline, leftmargin=2.8em, labelsep=0.5em, itemsep=0.35em]
		\item[\textbf{\textcolor{CTrap}{易错点二（内外混淆）}}]
		      将正方体内切球半径 $\frac{a}{2}$ 当作外接球半径。

		\item[\textbf{\textcolor{CTrap}{正确理解（半径区分）：}}]
		      外接球过正方体的顶点，半径是正方体中心到顶点的距离：
		      \[
			      R=\frac{\sqrt{3}}{2}a.
		      \]
		      内切球与正方体的面相切，半径是正方体中心到面的距离：
		      \[
			      r=\frac{a}{2}.
		      \]
		      二者不同。

		\item[\textbf{\textcolor{CTrap}{错因分析（距离对象混淆）：}}]
		      正方体的外接球与内切球球心都在正方体中心，并不是球心位置不同。

		      真正的区别在于半径对应的距离不同：外接球半径是“中心到顶点”的距离，内切球半径是“中心到面”的距离。
	\end{description}

	\medskip
	\begin{description}[style=nextline, leftmargin=2.8em, labelsep=0.5em, itemsep=0.35em]
		\item[\textbf{\textcolor{CTrap}{易错点三（逆解失误）}}]
		      已知外接球半径 $R$ 反求棱长时，错误得到
		      \[
			      a=\frac{\sqrt{3}}{2}R
		      \]
		      或
		      \[
			      a=\frac{R}{\sqrt{3}}
		      \]
		      等。

		\item[\textbf{\textcolor{CTrap}{正确理解（公式变形）：}}]
		      由
		      \[
			      R=\frac{\sqrt{3}}{2}a
		      \]
		      正确变形得
		      \[
			      a=\frac{2R}{\sqrt{3}}=\frac{2\sqrt{3}}{3}R.
		      \]

		\item[\textbf{\textcolor{CTrap}{错因分析（代数错误）：}}]
		      解方程时应两边同乘 $\frac{2}{\sqrt{3}}$，不要把乘除关系弄反，也不要误写成带根号的形式，例如 $\sqrt{\frac{2R}{3}}$。
	\end{description}
\end{trapbox}